\subsection{Higher connectivity and the finite-level comparison}

Connectivity of a simplicial profinite set below means connectivity
in Quick's homotopy theory: after fibrant replacement its component
set is a point and its profinite homotopy groups vanish in the stated
positive degrees.  Exact connectivity $c\geq0$ means also that
$\Pi_{c+1}$ is nonzero.  Exact connectivity $-1$ means nonempty and
disconnected.  All finite coefficient groups in this subsection are
arbitrary finite abelian groups with constant action unless an action
is explicitly specified.

\begin{result}{D.35}{Proved}{Higher finite-level connectivity in genera zero, one and two}
Let $Y_k=P_k^{\rm lev}$ denote the full level completion and put
$d=3g-3+n$.  For every stable type of genus at most two, the natural
map $Y_k\to\check P_k$ is the canonical isomorphism of D.19.
For $k\geq d$ these profinite spaces are contractible.  For
$0\leq k<d$ the following assertions hold.

In genus zero or one their connectivity is exactly $k-1$.
In genus two their connectivity is at least $-1,0,1$ at
$k=0,1,2$, respectively, and at least $k-2$ for $k\geq3$.
In particular the full nonempty curve complex has profinite
connectivity at least $h(g,n)-1$ in all these genera.
The requested four low types have the following \emph{exact}
connectivities, where $\ast$ denotes a contractible cone:
\[
\begin{array}{c|ccc}
(g,n)&P_1&P_2&P_3\\\hline
(1,2)&0&\ast&\ast\\
(1,3)&0&1&\ast\\
(2,0)&0&1&\ast\\
(2,1)&0&1&1.
\end{array}                                                    \tag{D.61}
\]
Also $P_3(2,2)$ has exact profinite connectivity one.  Thus both
counterexamples to the proposed connectivity guesses, in D.6 and
D.7, persist in the full and procongruence completions.

There is an exact further comparison in arbitrary genus.  For
$2\leq k\leq6$ and $k<d$, put
\[
c(g,n,k)=
\begin{cases}
k-1,&g=0,\\
\displaystyle\max\left\{k-g,
 \min\left(h(g,n),\left\lfloor\frac{2(k+1)}3\right\rfloor\right)-1
 \right\},&g\geq1.
\end{cases}                                                     \tag{D.62}
\]
For every finite abelian group $A$ and $0\leq j\leq c(g,n,k)$,
the actual Borel and orbit maps induce natural isomorphisms
\[
 H^j_{\rm cts}(Y_k;A)
       \ \cong\ \varinjlim_{H\triangleleft_f\Gamma_{g,n}}
                         H^j(H;A).                              \tag{D.63}
\]
The maps in the direct limit are restrictions.  Formula (D.63)
identifies the remaining full-level cohomology groups in this
range; it asserts no vanishing of its right side in higher genus.
The procongruence fundamental groups in every rank remain the
exact quotients of D.34.
\end{result}
\begin{proof}
We first establish the group and topology input.  Free groups,
closed orientable surface groups, and free abelian groups are good.
The surface statement and the permanence under finite-index
passage, products and the extensions used here are
\cite[Lemmas 3.1--3.3, Proposition 3.4, and Proposition 3.6]{D_GJZ2008}.
The opened surface proof uses the separating-circle amalgam,
the full induced topology on its free factors and cyclic edge
group, and the cohomology Mayer--Vietoris comparison.  These are
the precise goodness properties needed here.

Successive ordered-configuration fibrations prove goodness of
the genus-zero pure groups.  Their free kernels are finitely
generated and have finite cohomology with finite coefficients;
the actual completed extensions are exact, so comparison of the
two Hochschild--Serre spectral sequences applies.  Finite
extensions give the mixed groups as well.
$\Gamma_{1,1}=\mathrm{SL}_2(\mathbb Z)$ is virtually free.
A finite-index subgroup of the closed genus-two group avoiding
its central involution embeds with finite index into the
six-point sphere group.  Such a subgroup exists by residual
finiteness, already supplied by D.19.  Hence the closed genus-two
group is good.  Puncturing uses the actual completed Birman
extensions of D.31, with a surface kernel for the first closed
puncture and free kernels thereafter.  The same spectral-sequence
argument proves goodness for all $\Gamma_{g,n}$ with $g\leq2$.

An actual multicurve stabilizer has its pinned central cutting
extension by a free abelian twist lattice.  Its quotient is the
product of the cut-component groups.  In genus at most two the
comparison in D.32, using D.19 and the valid Boggi cutting theorem,
proves both exactness on completion and the full topology induced
from the ambient group.  The quotient and lattice are good and
have finite cohomology with finite coefficients; the same two
spectral sequences prove goodness of the pinned stabilizer.
The actual finite permutation extension preserves goodness and
full topology.  Stabilizers of chains are finite-index subgroups
of these groups and have both properties too.  For arbitrary
ambient genus, $k\leq6$ and $k<d$, the cut-component dimension
argument of D.32 gives genus at most two for every component of
every vertex of $P_k$.  Consequently all the same local
stabilizer statements hold in the precise range of (D.63),
without assuming goodness of the ambient group.

Here is the comparison on spaces.  Fix one contractible free
$\Gamma=\Gamma_{g,n}$-complex $E\Gamma$.  For every finite-index
normal subgroup $H$, the two actual maps are
\[
 P_k/H\ \xleftarrow{\ p_H\ }\ E\Gamma\times_H P_k
                  \ \xrightarrow{\ q_H\ }\ E\Gamma/H=BH.
                                                               \tag{D.64}
\]
The first forgets the $E\Gamma$ coordinate; the second forgets
the $P_k$ coordinate.  They commute with every finite-level
restriction.  Work on the nerve of the multicurve poset, so
cell stabilizers fix cells pointwise.  Filtering the middle space
by the dimension of its $P_k$ coordinate gives the cellular
equivariant spectral sequence
\[
 E_1^{p,q}(H)=
   \prod_{[s]\in H\backslash(P_k)_p}H^q(H_s;A)
      \ \Longrightarrow\ H^{p+q}(E\Gamma\times_H P_k;A).
                                                               \tag{D.65}
\]
It is obtained by replacing each orbit $p$-cell by its stabilizer
classifying space times that cell.  This construction also shows
that the row $q=0$ is the cellular cochain complex of $P_k/H$
and that its edge map is exactly $p_H^*$.  There are finitely
many cell orbits in every degree and the cell dimension is
bounded independently of $H$.

Every positive-degree cohomology class of a good group with
finite coefficients dies on a subgroup of finite index.
Indeed goodness represents it by a continuous profinite
cocycle, which factors through a finite quotient; on the
kernel a normalized positive-degree cocycle is zero.
Finite-index subgroups are good as well.
For a fixed $H,p,q>0$, an element of the product in (D.65)
is a finite tuple.  Kill each component on a finite-index
subgroup of its stabilizer.  Full ambient topology permits
one finite-index subgroup of $H$ whose intersections with all
these stabilizers lie in the chosen subgroups.  Replace it by
its normal core in $\Gamma$.  This remains inside $H$ and
inside the stated normal-level indexing system.  Every new orbit at that smaller level
then has a stabilizer conjugate to a subgroup on which one of
the original classes was killed.  Thus the whole tuple
restricts to zero.  This proves
$\varinjlim_H E_1^{p,q}(H)=0$ for all $q>0$; it does not
assume that only finitely many orbits occur over all levels.

For $q=0$, continuous functions to the finite group $A$ factor
through a finite level, so the colimit row is exactly
$C^*_{\rm cts}(Y_k;A)$.  Filtered colimits are exact, and the
bounded cell filtration gives a finite filtration in each
total degree.  Taking colimits in (D.65) therefore shows that
$\varinjlim p_H^*$ is an isomorphism in every degree.
The second map in (D.64) is the Borel fibration with fibre
$P_k$.  The discrete connectivity estimate proved in D.36
below and its cohomology spectral sequence show that $q_H^*$
is an isomorphism through degree $c(g,n,k)$.  Hence (D.63)
is the composite of $(\varinjlim q_H^*)^{-1}$ and
$\varinjlim p_H^*$; this specifies both comparison maps.
In genus at most two the same argument applies in every
rank $k<d$, with each of the stronger discrete connectivity
bounds of D.36.

For these low genera, ambient goodness kills
$\varinjlim_H H^j(H;A)$ for $j>0$ by the same restriction
argument.  Thus $H^j_{\rm cts}(Y_k;A)=0$ in every positive
degree up to the corresponding discrete connectivity bound,
for every finite $A$.  D.11 gives $\Pi_1(Y_k)=1$ when
$k\geq2$.  To pass to higher homotopy, take a fibrant
replacement.  Its finite-coefficient cohomology is unchanged.
Continuous cochains with values in $\mathbb Q/\mathbb Z$
are the union of the cochains with finite cyclic values.
Pontryagin duality identifies their cohomology with the
character group of the profinite integral homology: duality
is exact on the compact chain groups and their closed
boundary subgroups.  The finite-coefficient vanishing thus
implies vanishing of that homology in the stated degrees.
Induction on the first possibly nonzero homotopy group,
using Quick's Hurewicz theorem for degrees at least two
\cite[Theorem 2.30, p.~602]{D_Quick2008}, proves the claimed
higher connectivity.  Nonemptiness and connectedness in
degrees $-1,0$ follow from the finite-level quotients and
compactness, as in D.11.  The empty multicurve gives a
compatible cone contraction at all levels when $k\geq d$.

We finally prove the nonvanishing assertions, rather than
infer them from the discrete homotopy type alone.  The actual
curve complex embeds into its procongruence completion by
the audited Boggi Proposition 3.3 and Theorem 4.2.  Hence any
finite set of actual curves is separated at one finite
congruence level.  The same is true of the finite multicurves
and flags formed from that set: equality of their quotient
flags would identify their finite sets of curve vertices.
Thus a finite cellular cycle can be sent injectively on its
support into a sufficiently fine level quotient.

In genera zero and one, D.36 supplies a nonzero finite
$k$-cycle in $P_k$ with $\mathbb F_2$ coefficients, and
$P_k$ has dimension $k$.  Its separated image at a finite
level is a nonzero top-dimensional cycle.  A cochain dual
to a cell occurring with coefficient one is a cocycle,
since there are no higher cells; its evaluation is one.
Pull it back to $Y_k$.  Its value on the original actual
cycle remains one, so it cannot become a continuous
coboundary.  Together with the proved $(k-1)$-connectivity
this gives $\Pi_k(Y_k)\ne0$ for $k\geq1$; for $k=0$ it
separates two components.  This proves exactness in those
genera.  The same top-cell argument applies to the actual
graphs $P_1$ and two-dimensional complexes $P_2$ in D.6,
proving every indicated $0$ and $1$ in (D.61) except its
last entry.

For that entry, ordinary puncture forgetting gives a
continuous map of full level complexes
$\mathcal C_{2,1}^{\rm lev}\to\mathcal C_{2,0}^{\rm lev}$.
Continuity follows from the actual forgetting homomorphism:
the preimage of each finite-index subgroup of
$\Gamma_{2,0}$ has finite index in $\Gamma_{2,1}$.
Choose the nonzero finite two-cycle in
$\mathcal C_{2,0}=P_2(2,0)$ and its detecting finite-level
cochain just constructed.  Represent its finitely many
curves by geodesics in one fixed metric.  Choose a point
outside their union and move it to the prescribed puncture
by an isotopy.  The image curves give a finite two-cycle
in $\mathcal C_{2,1}$ mapping to the original one under
forgetting.  Pulling back the detecting cochain therefore
gives a nonzero continuous class in degree two.  This
proves the last entry since D.11 already gives simple
connectivity.  Explicitly, a nonzero such
$\mathbb F_2$ class injects into
$H^2(-;\mathbb Q/\mathbb Z)$ by the coefficient sequence
$0\to\mathbb F_2\to\mathbb Q/\mathbb Z
 \xrightarrow{2}\mathbb Q/\mathbb Z\to0$ and the
vanishing of $H^1(-;\mathbb Q/\mathbb Z)$.
Pontryagin duality and Hurewicz then give
$\Pi_2\ne0$ without a universal-coefficient assumption.

For $P_3(2,2)$ use the actual forgetting retraction to
$\mathcal C_{2,1}$ and the section of D.7.  The same finite
two-cycle has a closed lift through that section.  The
composite of the two actual forgetting maps is continuous
on finite levels and pulls back the preceding detecting
cochain to a nonzero class.  Its lower connectivity is one
by D.36.  This proves the asserted second completed
counterexample.  D.19 identifies all these full-level
statements with their procongruence counterparts.
\end{proof}

\begin{result}{D.36}{Proved}{A sharper rank bound and its exact upper-rank range}
For every stable positive-genus type and $0\leq k<d$,
the discrete complex $P_k$ is at least
$c(g,n,k)$-connected, with $c$ as in (D.62).
For genus zero the connectivity is exactly $k-1$.
For genus one it is also exactly $k-1$.
For every $g\geq1$ the further range
\[
                 3g-3\leq k\leq d-2                         \tag{D.66}
\]
has exact connectivity $k-g$, with nonzero reduced integral
homology in degree $k-g+1$.

In particular the complete discrete genus-two answer is:
$P_0$ is nonempty and disconnected; $P_1$ has exact
connectivity zero; $P_2$ has exact connectivity one;
and every $P_k$ with $3\leq k<d$ has exact connectivity
$k-2$.  At $k\geq d$ the complex is the contractible cone
already specified.  The assertions about $P_1$ and $P_2$
only apply when those ranks are strictly below $d$.
\end{result}
\begin{proof}
Let a nonempty multicurve have $q$ edges, dual graph first
Betti number $b$, and $v_+$ positive-genus and $v_0$ rational
vertices.  Since $g=\sum_v g_v+b>0$, one has
$b+v_+\geq1$.  The graph Euler equation gives
\[
       q-v_0=b+v_+-1\geq0,\qquad
       H(\sigma)=2g-3+n-v_0\geq r-g,
       \quad r=d-q.                                          \tag{D.67}
\]
The second expression is the exact upper-link dimension
from D.5, not the dimension of the original multicurve.
When deleting barycentric vertices to obtain $P_k$, every
deleted vertex has $r\geq k+1$.  Its upper link is therefore
at least $(k-g)$-connected.  The same coning argument as
in D.5 proves the additional lower bound $k-g$.
The initial Harer connectivity is at least this number:
for $n>0$ it follows from $k\leq3g-4+n$ and
$h-1=2g-4+n$; for closed stable types,
$h-1=2g-3$ gives the same inequality.
Taking the larger of this estimate and D.5 proves (D.62).

In genus one $b+v_+=1$, so $v_0=q$ and the upper-link
dimension is exactly $r-1$.  The deletion proof gives
$(k-1)$-connectivity, just as in genus zero.  These
complexes have dimension $k$.  To exhibit a nonzero
top-dimensional cycle, fix a multicurve $\tau$ leaving
one component $S_{0,k+4}$ in genus zero, or $S_{1,k+1}$
in genus one, with every other component a pair of
pants.  Such a choice has $|\tau|=d-k-1$.  Explicitly,
when $\tau$ is nonempty, put the excess labelled points
on one rational tail and decompose that tail into pants;
the core has the specified number of remaining flags.
The map of face posets
$\rho\mapsto\tau\cup\rho$ includes the curve complex of
that core into $P_k$.  Harer's sphere theorem supplies
a nonzero finite $k$-cycle on the core, also over
$\mathbb F_2$.  Its image is a nonzero chain in the
$k$-dimensional complex $P_k$, so is not a boundary.
For $k=0$ take the difference of two distinct vertices.
This proves exact connectivity in genera zero and one,
including the nonvanishing input used in D.35.

For (D.66), set $n'=k-3g+5$ and $N=n-n'$.  The two
inequalities say exactly that $n'\geq2$ and $N\geq0$.
Forgetting those $N$ labels gives an order map
\[
 r:P_k(g,n)\longrightarrow
       P_k(g,n')=P_{d'-2}(g,n'),\quad d'=3g-3+n'=k+2.
                                                               \tag{D.68}
\]
Here is the edge-count check needed for this map.
On forgetting one marked point, the only potentially
unstable vertex is the one carrying that point.  If it
becomes unstable, it is rational with two remaining
flags.  A tail with one leg and one edge loses that
edge; a component with two incident edges merges those
edges.  In either case the number of essential
multicurve edges decreases by one.  A neighboring
vertex exchanges one flag for one flag, so there is
no further cascade.  The exceptional entire unstable
curve does not occur here because $g\geq1$ and the
target has at least two marked points.  Thus after
$N$ forgettings the size decreases by at most $N$.
A source multicurve has at least
$d-k=N+2$ edges, so its image has at least two,
as required in the target.  Forgetting preserves
inclusions after inessential curves and duplicate
isotopy classes are removed, proving the order-map
assertion with its exact stabilization convention.

The map has an explicit section.  Choose one retained
label and replace its punctured collar by a rational
tail carrying that label and the $N$ new labels.
Choose $N$ nested pants cuts on that tail.  The fixed
collar identification lifts every target multicurve
to the complementary core; adjoining the $N$ cuts
defines an order map $i$ into $P_k(g,n)$.  It adds
exactly $N$ edges, and forgetting contracts precisely
that tail.  Hence $r i$ is the identity on every
multicurve and every inclusion.  For $N=0$ these are
the identity maps.  Realization gives a genuine
retraction, not only a homology comparison.

D.7 applies to the target of (D.68).  Its sphere
dimension parameter is $h'=2g-3+n'=k-g+2$, and it
has a split nonzero reduced homology group in degree
$h'-1=k-g+1$.  Applying $i_*$ and $r_*$ retains this
nonzero group in $P_k(g,n)$.  The lower bound
$k-g$ from (D.67) is therefore exact.  This also
explains why increasing the number of forgotten
labels requires the edge-count argument rather than
repeating a claim that every individual curve survives.

For genus two, (D.66) covers $3\leq k\leq d-2$.
At the remaining full rank $k=d-1$ Harer's theorem
gives exact connectivity $h-1=k-2$ when $n>0$.
For the closed surface $d=3$ the remaining rank is
$k=2$ and the connectivity is one, as in D.6.
The lower bounds at ranks one and two are D.5.
Their nonzero top homology for every $n$ can be seen
by fixing complementary pants cuts.  For $P_1$ leave
two once-holed tori and pants; their joined curve
complexes contain a square.  With marked points,
the rational pants between the two tori carry those
labels.  For $P_2$ and $n\geq1$, first cut a
nonseparating curve and put all the ordinary marked
points on a decomposed rational tail when necessary;
the remaining core is $S_{1,3}$, whose curve complex
has nonzero two-dimensional homology.  Its inclusion
by union with the fixed cuts lies in the
two-dimensional $P_2$.  For $n=0$, use the full
$\mathcal C_{2,0}$ of D.6.  These cycles cannot bound
in the respective dimensions.  Distinct pants
decompositions give at least two vertices of $P_0$,
which is a discrete complex.  This proves all the
stated genus-two ranks and completes the proof.
\end{proof}
